\documentclass{article}
\usepackage[UTF8]{ctex}
\usepackage{amsmath,amssymb,amsthm}
\usepackage{algorithm,algorithmicx,algpseudocode}
\usepackage{graphicx,booktabs,multirow}

\begin{document}
在机器学习中，我们可能会思考这些问题：学习意味着什么？机器学习算法为什么有效？什么样的问题是机器可以学习的？PAC learning（Probably Approximately Correct learning）就是一个可以回答上述问题的框架，它试图通过数学推理和证明为机器学习奠定理论基础。PAC learning由 Leslie Valiant教授于1984 年发明，它催生了计算机科学的一个新子领域，称为计算学习理论。

\section{PAC学习理论}
\subsection{一个例子}
想象一下两个玩家之间的猜区间游戏。玩家1的脑海里有一个区间[a,b]，他通过某种固定的方式随机生成一些数字，他不断向玩家2报告这些数字，并说明这些数字是否在区间内。玩家2希望通过玩家1的报告的数字以及标签猜测出玩家1脑海里的那个区间。玩家2得到的信息是一堆样本，每个样本包含一个数字和相应的标签，他的目标是确定区间断点a和b。\\
如果端点可以是任意实数，则玩家2几乎无法准确猜测区间，因为玩家1给出的样本数量有限。但是，无论玩家2在最后猜测的区间是什么，都可以根据玩家1的数字生成方案进行测试。于是说，我们可以验证玩家2答案的准确性。在这个猜数游戏中，如果玩家2能以很高的概率，返回一个近似正确的答案区间，那么这个问题就是PAC-learnable（PAC，即Probably Approximately Correct，概率近似正确）。\\
那么我们如何数学的刻画这个问题呢？
\subsection{一些基本概念}
我们首先介绍一些基本概念，首先我们有一个未知但是固定的分布D，数据样本都是独立的从分布D中产生。
\begin{definition}[概念类/目标类（Concept Class C）]
一组待学习的函数集合（比如刚才的区间）。
\end{definition}
\begin{definition}[假设类（Hypothesis Class H）]
学习算法可能输出的假设集合
\end{definition}
\begin{definition}[误差]
输出的假设和概念在分布D下不同的概率
\end{definition}
输出的假设c和概念h在数据分布D下的误差记为
\[ \text{erre}, p(h) = P_p(h(x) \neq c(x)) \]

\begin{definition}
    在集合 \( X \) 上的概念类 \( C \) 被称为可高效 \textbf{PAC} 学习（使用假设类 \( H \)），如果存在一个算法 \( A(\epsilon, \delta) \)，该算法能够访问 \( C \) 的查询函数，且其运行时间为 \( O(\text{poly}(1/\epsilon, 1/\delta)) \)，使得对于所有 \( c \in C \)、所有 \( X \) 上的分布 \( D \)，以及所有 \( 0 < \delta, \epsilon < 1/2 \)，算法 \( A \) 能以至少 \( 1 - \delta \) 的概率生成一个假设 \( h \in H \)，且该假设的误差不超过 \(\epsilon\)。

    \[
    P_D(\text{err}_{c,D}(h) \leq \epsilon) \geq 1 - \delta
    \]
\end{definition}
 该定义表达的核心意思是，在多项式时间内，是否存在算法通过数据集能以高概率学习到一个近似正确的假设，如果可以，该问题就是PAC-Learnable。\\
 根据我们刚才的定义，区间学习的例子是PAC-Learnable，大致证明思路：首先是一个针对区间预测的朴素的算法：在数据集中找到最大的正例和最小的正例，将他们作为我们得出的区间的边界。然后只需计算按照这种算法，得到的区间的错误概率，得到样本数量和误差的关系式，控制样本数量足够大，我们可以以大概率保证我们的误差在一个小的范围内。

\subsection{PAC学习理论}


\subsection{基本定义}
\begin{definition}[PAC学习]
对于概念类$\mathcal{C} \subseteq \{0,1\}^\mathcal{X}$，若存在算法$\mathcal{A}$和多项式函数$poly(\cdot,\cdot,\cdot,\cdot)$，使得$\forall \epsilon,\delta \in (0,1)$，$\forall c \in \mathcal{C}$，当样本量$m \geq poly(1/\epsilon,1/\delta,\text{size}(c),n)$时，有：

\begin{equation}
\Pr_{S\sim\mathcal{D}^m}\left[R(\mathcal{A}(S)) \leq \min_{h\in\mathcal{H}}R(h) + \epsilon\right] \geq 1-\delta
\end{equation}

则称$\mathcal{C}$是PAC可学习的。
\end{definition}

\subsection{区间预测问题分析}
考虑区间概念类$\mathcal{C} = \{\mathbb{I}_{[a,b]}: a,b \in \mathbb{R}\}$，其对应的经验风险最小化(ERM)算法如算法\ref{alg:interval}所示。

\begin{algorithm}[t]
\caption{区间ERM算法}
\label{alg:interval}
\begin{algorithmic}[1]
\Require 训练样本$S = \{(x_i,y_i)\}_{i=1}^m$，其中$y_i = \mathbb{I}_{[a^*,b^*]}(x_i)$
\Ensure 预测区间$[\hat{a},\hat{b}]$
\State 提取所有正例样本$S^+ = \{x_i | y_i = 1\}$
\State 计算$\hat{a} = \min S^+$, $\hat{b} = \max S^+$
\State \Return $[\hat{a},\hat{b}]$
\end{algorithmic}
\end{algorithm}

\begin{theorem}
\label{thm:interval}
对于任意$\epsilon,\delta > 0$，当样本量满足：

\begin{equation}
m \geq \frac{2}{\epsilon}\ln\frac{2}{\delta}
\end{equation}

时，算法\ref{alg:interval}输出的假设$\hat{h}$满足$R(\hat{h}) \leq \epsilon$的概率至少为$1-\delta$。
\end{theorem}

\begin{proof}
定义边界误差事件：
\begin{align*}
E_1 &= \{\exists x \in [a^*,b^*]: x < \hat{a}\} \\
E_2 &= \{\exists x \in [a^*,b^*]: x > \hat{b}\}
\end{align*}

由Chernoff不等式可得：
\begin{equation}
\Pr(E_1) \leq e^{-m\epsilon/2}, \quad \Pr(E_2) \leq e^{-m\epsilon/2}
\end{equation}

通过联合界(union bound)得到总误差概率：
\begin{equation}
\Pr(R(\hat{h}) > \epsilon) \leq \Pr(E_1) + \Pr(E_2) \leq 2e^{-m\epsilon/2} \leq \delta
\end{equation}
\end{proof}

\section{VC维理论}
\label{sec:vc}

\subsection{定义与性质}
\begin{definition}[VC维]
假设类$\mathcal{H}$的VC维定义为：
\begin{equation}
\text{VCdim}(\mathcal{H}) = \sup\{m: \Pi_\mathcal{H}(m) = 2^m\}
\end{equation}
其中$\Pi_\mathcal{H}(m)$为增长函数。
\end{definition}

\begin{theorem}[Sauer引理]
对于任意$\mathcal{H}$，当$m \geq d$时：
\begin{equation}
\Pi_\mathcal{H}(m) \leq \sum_{i=0}^d \binom{m}{i} \leq \left(\frac{em}{d}\right)^d
\end{equation}
\end{theorem}

\subsection{泛化误差界}
\begin{theorem}
对于VC维为$d$的假设类$\mathcal{H}$，以概率至少$1-\delta$有：
\begin{equation}
\sup_{h\in\mathcal{H}} |R(h) - \hat{R}(h)| \leq \sqrt{\frac{2d\log(em/d)}{m}} + \sqrt{\frac{\log(1/\delta)}{2m}}
\end{equation}
\end{theorem}

\section{在SVM中的应用}
\label{sec:svm}

对于线性分类器$\mathcal{H} = \{x \mapsto \text{sign}(w^\top x + b): \|w\| \leq \Lambda\}$，其VC维满足：

\begin{equation}
d_{\text{VC}} \leq \min\left(\lceil R^2\Lambda^2 \rceil + 1, n+1\right)
\end{equation}

其中$R = \sup\|x\|$为数据半径。通过核函数$\phi$映射后，有效VC维由间隔$\gamma$控制：

\begin{equation}
d_{\text{eff}} \approx \left\lceil \frac{R_\phi^2}{\gamma^2} \right\rceil
\end{equation}

\section{深度学习中的理论挑战}
\label{sec:dl}

传统理论预测神经网络需要样本量：

\begin{equation}
m \geq \frac{C\cdot\text{VCdim}(\mathcal{H})}{\epsilon}
\end{equation}

但实际观察到的现象与理论不符：

\begin{table}[t]
\centering
\begin{tabular}{lcc}
\toprule
模型结构 & 理论所需样本量 & 实际有效样本量 \\
\midrule
ResNet-50 & $\sim 10^{10}$ & $\sim 10^6$ \\
Transformer & $\sim 10^{12}$ & $\sim 10^7$ \\
\bottomrule
\end{tabular}
\caption{VC维预测与实际样本需求的巨大差距}
\label{tab:gap}
\end{table}

\section{结论}
本文系统分析了PAC学习和VC维的理论体系，揭示了其在解释现代深度学习模型时的局限性。未来的研究方向包括：

\begin{itemize}
\item 发展新的容量度量替代VC维
\item 研究梯度下降的隐式正则化效应
\item 探索神经网络的低维本质特性
\end{itemize}



\end{document}